%%
%% Test: Layout Configuration — v1.1
%% Stage: layout refactor (wide margin + sidenotesplus + textbook leading)
%% 
%% Purpose: Verify layout settings:
%%   - Page size (17.8 × 25.4 cm)
%%   - Asymmetric margins with wide outer margin for notes
%%   - Header/Footer (chapter name / page number)
%%   - Line spacing (1.15 — Boyer/Stewart-style)
%%   - Paragraph indent (1em)
%%   - Margin notes via sidenotesplus (\mnote, \marginfig)
%%   - draftwatermark (conditional)
%% 

\documentclass{matinbook}

\title{تست چیدمان صفحه — نسخه ۱.۱}
\author{تیم MatinBook}
\date{۱۴۰۵}

\booktitle{تست چیدمان MatinBook}

\begin{document}

\maketitle

\tableofcontents

\chapter{بررسی چیدمان صفحه}
\label{chap:test-layout}

این سند، تنظیمات چیدمان \texttt{mb-layout} را بررسی می‌کند.
هدف این است که ببینیم آیا ابعاد صفحه، حاشیه‌های نامتقارن،
یادداشت‌های حاشیه، فاصله‌ی خطوط، و سرصفحه/پاصفحه به‌درستی
اعمال می‌شوند یا نه.

%----------------------------------------
\section{قطع صفحه}
\label{sec:page-size}

\textbf{عرض صفحه:} \the\paperwidth

\textbf{ارتفاع صفحه:} \the\paperheight

\textbf{عرض متن:} \the\textwidth

\textbf{ارتفاع متن:} \the\textheight

قطع صفحه‌ی فعلی \texttt{17.8 × 25.4} سانتی‌متر است که معادل
\texttt{7 × 10} اینچ می‌باشد و در کتاب‌های فنی آمریکایی
(مثل O'Reilly و Boyer) رایج است. این قطع، تعادل خوبی بین
قابلیت حمل و فضای کافی برای فرمول‌های ریاضی برقرار می‌کند.

%----------------------------------------
\section{حاشیه‌ها}
\label{sec:margins}

حاشیه‌های صفحه به‌صورت نامتقارن تنظیم شده‌اند تا فضای
بیشتری برای یادداشت‌های حاشیه فراهم شود:

\textbf{حاشیه‌ی بالا:} ۱.۷ سانتی‌متر

\textbf{حاشیه‌ی پایین:} ۱.۷ سانتی‌متر

\textbf{حاشیه‌ی داخلی (عطف):} ۱.۵ سانتی‌متر

\textbf{حاشیه‌ی بیرونی:} ۵.۶ سانتی‌متر

\textbf{عرض ناحیه‌ی یادداشت:} ۳.۵ سانتی‌متر

\textbf{فاصله‌ی یادداشت از متن:} ۰.۳ سانتی‌متر

\textbf{جبران عطف (binding offset):} ۰.۴ سانتی‌متر

این ساختار باعث می‌شود متن اصلی باریک‌تر شود، ولی فضای
کافی برای یادداشت‌های توضیحی، فرمول‌های کمکی، و ارجاعات
در حاشیه فراهم شود. \mnote{این یک یادداشت حاشیه‌ای است که در حاشیه‌ی بیرونی صفحه ظاهر می‌شود.} این سبک در کتاب‌های درسی
آمریکایی مثل بویس و استوارت بسیار رایج است.

%----------------------------------------
\section{سرصفحه و پاصفحه}
\label{sec:header-footer}

در صفحات زوج، نام فصل در سرصفحه‌ی سمت چپ نمایش داده می‌شود.

در صفحات فرد، نام فصل در سرصفحه‌ی سمت راست نمایش داده می‌شود.

شماره‌ی صفحه در پاصفحه، سمت بیرونی، قرار دارد.

\textbf{نکته:} در حالت \texttt{twoside}، سرصفحه‌ها به‌صورت
آینه‌ای تنظیم می‌شوند تا در صفحات زوج و فرد، خوانایی حفظ شود.
\mnote{سرصفحه‌ها در حالت دوطرفه به‌صورت آینه‌ای تنظیم می‌شوند.} این استاندارد در کتاب‌های چاپی رعایت می‌شود.

%----------------------------------------
\section{فاصله‌ی خطوط}
\label{sec:line-spacing}

\textbf{فاصله‌ی خطوط:} ۱.۱۵ (سبک کتاب درسی)

\textbf{تورفتگی پاراگراف:} ۱em

\textbf{فاصله‌ی پاراگراف:} ۰pt

فاصله‌ی خطوط \texttt{1.15} برای کتاب‌های درسی ریاضی مناسب است،
زیرا فضای کافی برای فرمول‌های داخل متن و فرمول‌های جدا فراهم
می‌کند و خوانایی را افزایش می‌دهد. در مقایسه، سبک O'Reilly
از فاصله‌ی کمتری (حدود ۱.۰۵) استفاده می‌کند، ولی کتاب‌های
درسی مثل بویس و استوارت از فاصله‌ی بیشتری بهره می‌برند.

این یک پاراگراف نمونه است که برای بررسی فاصله‌ی خطوط نوشته
شده است. در تایپوگرافی فارسی، فاصله‌ی خطوط معمولاً بین
۱.۰ تا ۱.۲ است. انتخاب مقدار مناسب، بستگی به نوع محتوا
و مخاطب دارد. برای کتاب‌های درسی، مقادیر بالاتر (۱.۱ تا
۱.۲) توصیه می‌شود، ولی برای کتاب‌های مرجع، مقادیر پایین‌تر
(۱.۰ تا ۱.۱) کافی است.

%----------------------------------------
\section{یادداشت‌های حاشیه}
\label{sec:margin-notes}

یادداشت‌های حاشیه یکی از ویژگی‌های مهم کتاب‌های درسی
آمریکایی هستند. در این بخش، چند نوع یادداشت را آزمایش
می‌کنیم.

\mnote{یادداشت اول: ساده.}
این متن اصلی است و در کنار آن، یک یادداشت ساده قرار دارد.

\mnote{یادداشت دوم: کمی طولانی‌تر تا ببینیم در چند خط شکسته می‌شود و آیا با متن اصلی تداخل می‌کند.}
ادامه‌ی متن اصلی. یادداشت‌های حاشیه معمولاً برای توضیح
یک اصطلاح فنی، ارجاع به منبع، یا نکته‌ی جانبی استفاده
می‌شوند.

\mnote{یادداشت سوم.}
متن بیشتر... متن بیشتر... متن بیشتر... متن بیشتر...
متن بیشتر... متن بیشتر... متن بیشتر... متن بیشتر...

یک تصویر نمونه هم در حاشیه قرار می‌دهیم:
\marginfig{example-image}{نمودار نمونه‌ی حاشیه}
این تصویر در حاشیه‌ی بیرونی صفحه ظاهر می‌شود و می‌تواند
برای نشان دادن نمودارهای کوچک، شکل‌های هندسی، یا
تصاویر کمکی استفاده شود.

%----------------------------------------
\section{محیط‌های علمی}
\label{sec:environments}

محیط‌های علمی باید با چیدمان جدید (حاشیه‌ی عریض و فاصله‌ی
خطوط ۱.۱۵) به‌درستی کار کنند.

\begin{definition}[عدد اول]
    عدد طبیعی بزرگ‌تر از ۱ که تنها دو مقسوم‌علیه مثبت دارد،
    \emph{عدد اول} نامیده می‌شود.
    \label{def:layout-prime}
\end{definition}

\begin{theorem}[قضیه فیثاغورس]
    در هر مثلث قائم‌الزاویه، مربع وتر برابر است با مجموع
    مربعات دو ضلع دیگر:
    \[
    a^2 + b^2 = c^2
    \]
    \label{thm:layout-pythagoras}
\end{theorem}

\mnote{این قضیه یکی از قدیمی‌ترین قضایای ریاضی است.}

\begin{proof}
    اثبات این قضیه با روش‌های مختلفی امکان‌پذیر است.
    یکی از ساده‌ترین روش‌ها استفاده از چهار مثلث قائم‌الزاویه
    است که در یک مربع بزرگ چیده می‌شوند.
    \label{prf:layout-pythagoras}
\end{proof}

\begin{example}[حل معادله درجه دوم]
    معادله $x^2 - 5x + 6 = 0$ را حل کنید.
    \label{ex:layout-quadratic}
\end{example}

\begin{solution}
    با تجزیه داریم:
    \[
    x^2 - 5x + 6 = (x - 2)(x - 3) = 0
    \]
    پس $x = 2$ یا $x = 3$.
\end{solution}

\begin{remark}
    فرمول عمومی $x = \frac{-b \pm \sqrt{b^2 - 4ac}}{2a}$
    هم جواب یکسان می‌دهد.
\end{remark}

\begin{exercise}
    معادله $x^2 + 2x + 1 = 0$ را حل کنید.
\end{exercise}

%----------------------------------------
\section{جدول}
\label{sec:table}

جدول‌ها باید در عرض جدید متن (که باریک‌تر شده) به‌درستی
جا شوند.

\begin{table}[h]
\centering
\caption{جدول با ستون‌های سفارشی}
\label{tab:layout-test}
\begin{tabular}{L{3cm} C{2cm} R{3cm}}
\toprule
\textbf{چپ} & \textbf{وسط} & \textbf{راست} \\
\midrule
متن & ۱۰ & ۲۰ \\
متن دوم & ۳۰ & ۴۰ \\
\bottomrule
\end{tabular}
\end{table}

%----------------------------------------
\section{کد}
\label{sec:code}

بلوک‌های کد هم باید در عرض جدید جا شوند:

\begin{latin}
\begin{minted}{python}
def greet(name):
    return f"Hello, {name}!"

print(greet("MatinBook"))
\end{minted}
\end{latin}

%----------------------------------------
\section{فرمول‌های طولانی}
\label{sec:long-formulas}

فرمول‌های طولانی باید در عرض باریک‌تر جدید به‌درستی شکسته
شوند یا در صفحه جا شوند:

\begin{align*}
    \int_{-\infty}^{+\infty} e^{-x^2} \, dx &= \sqrt{\pi} \\
    \sum_{n=1}^{\infty} \frac{1}{n^2} &= \frac{\pi^2}{6} \\
    \lim_{n \to \infty} \left(1 + \frac{1}{n}\right)^n &= e
\end{align*}

این فرمول‌ها در حالت عادی باید در عرض متن جا شوند. اگر
عرض متن کافی نباشد، ممکن است به بیرون سرریز کنند یا
با حاشیه‌ی یادداشت تداخل کنند.

%----------------------------------------
\section{نتیجه‌گیری}
\label{sec:conclusion}

اگر این سند بدون خطا کامپایل شود:

قطع صفحه (۱۷.۸×۲۵.۴) اعمال شده است.

حاشیه‌های نامتقارن با ناحیه‌ی عریض برای یادداشت اعمال شده‌اند.

یادداشت‌های حاشیه با \texttt{sidenotesplus} کار می‌کنند.

سرصفحه و پاصفحه به‌درستی کار می‌کنند.

شماره‌ی صفحه در پاصفحه‌ی سمت بیرونی قرار دارد.

فاصله‌ی خطوط ۱.۱۵ اعمال شده است.

تورفتگی پاراگراف ۱em است.

محیط‌های علمی با چیدمان جدید سازگارند.

\textbf{وضعیت تست:} قبول.

\end{document}